\section{他者用モジュールの詳細}\label{app:component}
他者用モジュールは，構成を4通りに変えても他者の好みを書き出さなかった（表\ref{tab:module}）．
事前登録した主の構成（自由出力モジュール，経験50\%条件）では，入れ替え応答 T2 と好み差応答 T5 はともに4組中1組でしか合格の形にならず，不合格であった．
合格の形になった組の値（T2 $+0.03$〜$+0.05$，T5 $+0.04$〜$+0.05$）は，本当の快・不快を入れた値（T2 $+0.60$，T5 $+0.42$）の 1/10 ほどで，組によって向きもそろわなかった．
出力 $Q2$ と本当の快・不快の相関は弱く負に寄り，$Q2$ の分散の 71〜86\% は，本当の快・不快が常に0である範囲 2〜5 が占めていた．
それでも，生成した $Q2$ を入れたときの他者側の特徴予測の誤差（組1〜4，緑好みデータで 0.116〜0.146，好み混合データで 0.121〜0.159）は，本当の快・不快を入れた場合（同じく 0.252〜0.276，0.241〜0.265）よりずっと小さかった．
% 出典：numbers_check.md（Claude Code 第2弾，保存値）。v12 まとめの 0.13〜0.18・0.25〜0.28 は数え方が違う。
モジュールは，予測に役立つ別の情報（範囲 2〜5 を通り道にした他者の位置）を伝えていた．
% 出典：v12 まとめ §3 表2，§4（正式の判定と注意 1〜4）。

この形は，経験80\%条件の段階Iの上で学習しても変わらなかった（探索．T2 は4組中2組，T5 は1組で合格の形．値は $+0.02$〜$+0.05$，$Q2$ と本当の快・不快の相関は4組とも負，分散の 52〜80\% が範囲 2〜5）．
出力を自分が経験した快・不快の形に限った形制限モジュール（T2 は4組中0組），学習データの視界で A-2 を区間ごとに隠した他者遮蔽学習（T2 は3組中0組で，大きく逆向き），出力を1つの符号 $u_t$ に絞った1ビットモジュールでも，好みを書き出す形は出なかった．
1ビットモジュールでは，$u_t$ の符号は正しい候補と偶然より高く一致した（経験50\%条件・経験80\%条件とも，事前登録した判定で3組中3組．経験50\%条件では $u_t$ の平均が好み×配置に沿って入れ替わったが，経験80\%条件では一致の度合いが弱く，BA は 0.51〜0.65）．
しかし，位置をそろえた緑好きと赤好きの A-2 で $u_t$ の違い（T5）はほぼ0で（経験80\%条件で $-0.006$〜$0.000$），好みと方角が食い違うエピソードで $u_t$ が好みに従う割合は経験80\%条件で 0.30〜0.41 と，方角の側に寄った（事前登録した基準で3組中0組が好み）．
1ビットに絞っても，モジュールが書いたのは A-2 の好みではなく，A-2 の進む方角（または位置）であった．
ただし，この結論は調べた4通りの構成，予測の誤差だけによる学習，この課題に限られる．
また，調べたのはモジュールの出力であり，モジュールや共有モジュールの内部状態に好みの情報がないことを示したものではない．
% 出典：v13 まとめ §4（2段目探索），v12 まとめ §3 表2（b 0/4，c 0/3），v14 まとめ §2 (4) 表，v14 ログ §28・§39。
% 1ビットの経験50% の (a)(b) は事前登録していない（v14 ログ §29）ので，本文では経験80% の値を主にした。

\subsection{他者用モジュールの各構成}
各構成の結果を表\ref{tab:module}にまとめる．

\begin{table}[tbp]
  \centering
  \small
  \caption{他者用モジュールの各構成の結果}
  \label{tab:module}
  % paper/tables/module.tex -- generated by analyze/paper/tables_paper.py; sources in paper/tables/README.md
\setlength{\tabcolsep}{3.5pt}%
\begin{tabular}{@{}llp{0.21\linewidth}ll@{}}
\toprule
\multicolumn{5}{@{}p{0.9\linewidth}@{}}{(a) 主の見分け（合格の組），合格の形の組の経験ありの値，$Q2$ と本当の快・不快の相関，範囲2〜5の分散の割合} \\
\midrule
モジュール・条件 & T2・T5 & 合格の形の組の値 & 相関 & 範囲2〜5 \\
\midrule
自由出力・経験50\% & 1/4・1/4 & T2 $+0.03$〜$+0.05$，T5 $+0.04$〜$+0.05$ & $-0.11$〜$-0.02$ & 0.71〜0.86 \\
自由出力・経験80\%（探索） & 2/4・1/4 & T2 $+0.02$〜$+0.05$，T5 $+0.02$〜$+0.04$ & $-0.165$〜$-0.003$ & 0.52〜0.80 \\
形制限・経験50\%（探索） & 0/4・1/4 & --- & $-0.16$〜$+0.12$ & 0 \\
他者遮蔽学習・経験50\%（探索） & 0/3・1/3 & T5 $+0.034$〜$+0.045$ & $-0.47$〜$+0.18$ & 0 \\
1ビット・経験50\% & ---・--- & T5 $-0.013$〜$-0.000$ & $+0.317$〜$+0.347$ & 0 \\
1ビット・経験80\% & ---・0/3 & T5 $-0.006$〜$-0.000$ & $+0.080$〜$+0.219$ & 0 \\
\midrule[\heavyrulewidth]
\multicolumn{5}{@{}p{0.9\linewidth}@{}}{(b) T3：他者側の特徴予測の誤差（生成／真値／零）と好み追従率} \\
\midrule
モジュール・条件 & 生成 & 真値 & 零 & 好み追従率 \\
\midrule
自由出力・経験50\% & 0.121〜0.159 & 0.241〜0.265 & 0.293〜0.349 & --- \\
自由出力・経験80\%（探索） & 0.112〜0.188 & 0.377〜0.414 & 0.298〜0.519 & --- \\
形制限・経験50\%（探索） & 0.161〜0.185 & 0.241〜0.265 & 0.293〜0.349 & --- \\
他者遮蔽学習・経験50\%（探索） & 0.125〜0.149 & 0.139〜0.165 & 0.133〜0.177 & --- \\
1ビット・経験50\% & 0.211〜0.215 & 0.248〜0.265 & 0.293〜0.304 & 0.273〜0.614 \\
1ビット・経験80\% & 0.205〜0.250 & 0.391〜0.414 & 0.298〜0.519 & 0.295〜0.409 \\
\midrule
\multicolumn{5}{@{}p{0.97\linewidth}@{}}{注：範囲2〜5の分散の割合は，形制限・他者遮蔽学習・1ビットでは作りの上で0．T3 は好み混合データ，組1〜4（1ビットは組1〜3，ep200）の範囲．他者遮蔽学習の T3 は記録にある値で，A-2 が隠れているフレームの自己の視界の誤差．1ビットの T5 の合格は事前登録した条件 (a)（経験50\%は事後の確認）．好み追従率は好みと方角が食い違う88エピソード，段階IIの最終5点の範囲（組1〜3）．} \\
\bottomrule
\end{tabular}

  % 表A2：4通りの構成（自由出力／形制限／他者遮蔽学習／1ビット）× 経験50%・経験80%。
  % 列：T2・T5 の合格の組，経験ありの値の範囲，本当の快・不快の値（T2 +0.60，T5 +0.42），Q2 と本当の快・不快の相関，
  %     範囲 2〜5 の分散の割合，T3 の生成と真値，1ビットの好み追従率。
  % 出典：v12 まとめ §3・§4，v13 まとめ §4，v14 まとめ §2 (4)，v14 ログ §28・§29・§39。
\end{table}

\paragraph{自由出力モジュール}
経験50\%条件の上で，事前登録した主の構成として学習した．
T2 は組1だけ，T5 は組2だけが合格の形になり，向きは組で逆になった（組2の T2，組1の T5 は逆向き）．
$Q2$ の分散の 71〜86\% が範囲 2〜5 で，値は $\pm1$ に張り付き，$Q2$ から A-1 自身の一歩の向きは読めなかった．
経験80\%条件の上では，G5a の関門に不合格だったため（自己位置の追跡も基準に届かなかった），探索として学習した．
T2 は組1・組3 で合格の形になったが，値は $+0.02$〜$+0.05$ で本当の快・不快を入れた値の 1/12〜1/30 にとどまり，組の差の大部分は経験なし条件の値が負であることから来ていた．
組4を加えた4組では，T2・T5 とも不合格であった．
% 出典：v12 まとめ §4，v13 まとめ §4（表2と注意 1〜4，s3 を加えた4組）。

\paragraph{形制限モジュール}
自由出力モジュールの結果を見た後に設計した探索である．
出力を，範囲 2〜5 を0とし，範囲 0・1 と 6・7 を逆の符号にした形に限った．
どちら側が快いかを決める符号は，4組中3組でほぼ一定のまま固定され，残る組2 では動いたが好みには沿わなかった．どの組でも，符号は好みや配置に応じて入れ替わらなかった．
モジュールは範囲ごとの強さで位置を伝えた．
T2 は4組中0組，T5 は1組で合格の形になった．
範囲 2〜5 だけを0にした参考の構成（組1のみ）では，判定の形は合格したが（T2 5点中5点，T5 5点中4点），値は本当の快・不快の 1/8〜1/20 で，相関は弱く負であり，好みを表す形とは読めなかった．
% 出典：v12 まとめ §3 表2，research_overview（2段目 b・案 A）。

\paragraph{他者遮蔽学習}
「A-2 が見えない間は，A-2 の好みを知らなければ A-2 の動きを予測できないので，モジュールは好みを表すようになる」という仮説を，学習の前に登録して確かめた（扱いは探索）．
形制限モジュールの形のまま，学習データの入力の視界から A-2 を 10〜30 ステップの区間で消した．
評価は，学習と同じく A-2 を隠したテストで，A-2 が隠れている場面について行った．
T2 は3組中0組で，経験ありの値は大きく負（逆向き）になり，T5 は1組で合格の形になった．
符号はほぼ一定で，$Q2$ と本当の快・不快の相関は弱いか逆向きであった．
% 出典：v12 まとめ §3 表2，v12 ログ §46・§50・§52。

\paragraph{1ビットモジュール}
出力をスカラー $u_t$ にし，$Q2_t=u_t\,c^{+}$ とした（事前登録）．
経験50\%条件の上では，$u_t$ の符号と正しい候補の一致の BA は 0.67〜0.73 で事前登録した判定に合格し，$u_t$ の平均は好み×配置のセルに沿って入れ替わった．
しかし，T5 は $-0.013$〜$0.000$（本当の快・不快で $+0.42$），T2 は $-0.019$〜$+0.041$（最終5点の範囲）で，同じ場所にいる A-2 の好みの違いを表していなかった．
好みと方角が食い違うエピソードで $u_t$ が好みに従う割合は，最終5点の範囲で 0.27〜0.61 で，偶然を上回る組はなかった．
ただし，経験50\%条件のこの比較と T5 の条件は事前登録しておらず，判定を見た後に方法を決めた確認である．
% 出典：Claude Code 第1弾（最終5点の範囲 0.273〜0.614，図5・表A2 と同じ数え方），v14 ログ §29（事前登録していない）。
%   v14 まとめの 0.33〜0.58 は数え方が違う（2026-10-26 決定で最終5点にそろえた）。
経験80\%条件の上では，一致の BA は 0.51〜0.65 と弱く，事前登録した二つの条件（T5 が正，食い違うエピソードで好みに従う）はどちらも3組中0組で満たされなかった．
% 出典：v14 ログ §28・§29・§39，v14 まとめ §2 (4)。
